\section{Game Logic and Server Authority}
\label{sec:authority}

\subsection{All state is 10 bytes owned by the STM32}
The whole game is one \code{game_t}: nine cells and a state byte, 10 bytes with no padding. The client never holds it; it receives
only checksummed text replies. The side to move is computed from the board (X when the counts are equal), so there is no separate turn variable that could
disagree with the board, and a client cannot make it disagree.

\subsection{A rejection cannot change state because every check precedes the first write}
Listing~\ref{lst:move} is \code{game_move}. Every error \code{return} appears before the line \code{g->board[index] = side}, and the
checks only read \code{*g}. The proof is by inspection: a path that returns an error has executed no write, so there is nothing to undo.
The other commands share the shape of one condition followed by at most one state change. The order of the checks (state, range, turn,
occupancy) is the order in Table~\ref{tab:errors}, and it is the same order the protocol promises, so the choice of NACK is deterministic.

\begin{figure}[t]
\begin{lstlisting}[float=false,language=C, caption={Validate-then-commit (\texttt{game.c}, comments trimmed).}, label=lst:move]
if (g->state != GAME_IN_PROGRESS) return GAME_ERR_NOGAME;
if (index < 0 || index >= GAME_CELLS)
                                  return GAME_ERR_RANGE;
const char side = game_side_to_move(g);
if (mark != 0 && mark != side)    return GAME_ERR_NOTURN;
if (g->board[index] != '.')       return GAME_ERR_OCCUPIED;
/* ---- commit: first write to *g ---- */
g->board[index] = side;
const game_result_t r = game_evaluate(g->board);
if (r != GAME_RESULT_NONE) g->state = GAME_OVER;
\end{lstlisting}
\end{figure}

\subsection{The AI opponent cannot bypass the rules}
In a \code{NEW:AI} game the server picks O's cell with \code{ai_pick}, which only reads a copy of the board, and then makes the move through the same \code{game_move} shown in Listing~\ref{lst:move}. It therefore cannot play an occupied or out-of-range cell. The AI flag is kept in the server structure, not in \code{game_t}, so the 10-byte game and the exhaustive core tests are unchanged. The search is minimax with alpha-beta pruning, so O never loses. We checked this on the PC by letting X try every possible sequence of moves: 681 games, 0 X wins (498 O wins, 183 draws), and 936 AI moves, none illegal (\tn{test_ai.c}, PC). A rejected \code{MOVE} does not trigger an AI move (\tn{test_rejected_moves_do_not_wake_the_ai}).

\subsection{Win and draw are decided by the server after each commit}
After a commit, \code{game_evaluate} tests the eight lines (three rows, three columns, two diagonals) before testing for a full board, so a move
that fills the last cell and completes a line is a win, not a draw. The server then ends the game and replies \code{WIN} or \code{DRAW}.
Table~\ref{tab:logic} lists the checks that back the argument. They run on the PC (game core alone), except the last row, which is the board suite.

\begin{table}[t]
\centering
\caption{Evidence for the authority argument (PC unless stated).}
\label{tab:logic}
\scriptsize
\begin{tabular}{@{}L{3.9cm}L{3.9cm}@{}}
\toprule
Check & Result \\
\midrule
Every error compared byte-for-byte (\code{memcmp} of the 10 bytes) & 0 failures in 50{,}546{,}417 checks \\
All positions reachable by legal play, each tried with every index and claim & 549{,}946 positions; 6{,}859{,}116 illegal attempts, all left the game identical \\
Random calls (fixed seed): six commands, indices $-3$..12 & 2{,}000{,}000 calls, invariants held; 22{,}948 wins, 3{,}003 draws \\
Independent Python brute force vs.\ C rules & 255{,}168 games: 131{,}184 X, 77{,}904 O, 46{,}080 draws; C reproduces them \\
All 8 lines, both players, over the serial link & \tn{test_each_line_wins}, 16 cases; passes on the board \\
All 936 AI moves, X plays every sequence & 681 games: 0 X wins, 0 illegal AI moves \\
\bottomrule
\end{tabular}
\end{table}

The exhaustive run covers the game core only, not framing, and it ran on a PC. It uses one fixed random seed, so it is repeatable
rather than broad. The link-level claim rests on the same invariants being tested through the serial port on the board
(Sec.~\ref{sec:valid}). Hostile clients are bounded the same way: a client that lies about the turn gets \code{NACK:NOTURN}, one that
resends a move gets \code{NACK:OCCUPIED} or \code{NACK:NOGAME}, and none of these changes the 10 bytes.
